\makeabstract
    {
    Investigating the Catalytic Activity of ONO- Group V Alkoxide Complexes for the Ring-Opening Polymerization of Biodegradable Materials
    }
    {
    Dennery\textsuperscript{Ender} Simon\textsuperscript{1}, Manuel Badena\textsuperscript{1}, Christopher B. Durr\textsuperscript{1,2}
    }
    {
    \textsuperscript{1} Department of Chemistry, Amherst College, Amherst, MA\\
\textsuperscript{2} The Ohio State University, Columbus, OH
    }
    {
    We are looking to expand the overlooked usage of Group V metals in the creation of catalysts that use ring-opening polymerization to synthesize biodegradable plastics.
    }
    {
    Using various commercially available compounds we synthesized a variety of unique ligands based on previous students' studies and relevant literature.  These ligands, produced under fume hoods in a normal atmosphere, were then transferred into an oxygen and water free glovebox and used to create catalysts.  We then observed the structures of these chemicals with Nuclear Magnetic Resonance (NMR) and Single Crystal X-Ray Diffraction (SCXRD) technology.  Once made, these catalysts were combined with various monomers in the glovebox and melted together at 140{\textdegree}C so their polymerization could be analyzed.
    }
    {
    We were able to observe the changes to the properties of catalysts when made with Niobium versus Tantalum as well as how different functional groups present on the ligands used affected our catalysts.  We crystalized some of these catalysts and, using the SCXRD, were able to view their real structures and how they compare to our NMR spectra.  By viewing the kinetics data from our polymerizations, we determined the different rates at which the catalysts are able to operate.
    }
    {
    By knowing how a variety of functional groups on our ligands changes the efficacy of our catalysts, we will be better able to isolate the most effective groups and in turn make better catalysts with these Group V metals.  As we make better catalysts, the biodegradable plastics we can produce will also improve and their uses in society can continue to expand.
    }

    \newpage
    